\documentclass[../ap_2013.tex]{subfiles}

\graphicspath{{./image/}}

\begin{document}

\setcounter{chapter}{1}
\chapter{対称行列}
\section{}
\(\lambda\)を固有値としたとき、固有値方程式より
\begin{equation}\begin{aligned}[b]
    0 &= \abs{A-\lambda I}\\
    &= \lambda^3 -\tr A \lambda^2 + \det A\tr A^{-1} \lambda - \det A\\
    &= \lambda^3 -7 \lambda^2 + 14 \lambda - 8\\
    &= (\lambda-1)(\lambda-2)(\lambda-4)\\
    \lambda &= 1,2,4
\end{aligned}\end{equation}
とわかる。

\(\lambda_1=1\)に対応する固有ベクトル\(v_1\)は
\begin{equation}\begin{aligned}[b]
    \begin{pmatrix}
        1 & -1 & 0\\
        -1 & 2 & -1\\
        0 & -1 & 1
    \end{pmatrix}v_1&=0\\
    \rightarrow\, v_1 &=\frac{1}{\sqrt{3}} \begin{pmatrix}
        1\\1\\1
    \end{pmatrix}
\end{aligned}\end{equation}

\(\lambda_2=2\)に対応する固有ベクトル\(v_2\)は
\begin{equation}\begin{aligned}[b]
    \begin{pmatrix}
        0 & -1 & 0\\
        -1 & 1 & -1\\
        0 & -1 & 0
    \end{pmatrix}v_2&=0\\
    \rightarrow\, v_2 &=\frac{1}{\sqrt{2}} \begin{pmatrix}
        1\\0\\-1
    \end{pmatrix}
\end{aligned}\end{equation}

\(\lambda_1=3\)に対応する固有ベクトル\(v_3\)は
\begin{equation}\begin{aligned}[b]
    \begin{pmatrix}
        -2 & -1 & 0\\
        -1 & -1 & -1\\
        0 & -1 & -2
    \end{pmatrix}v_1&=0\\
    \rightarrow\, v_3 &=\frac{1}{\sqrt{6}} \begin{pmatrix}
        1\\-2\\1
    \end{pmatrix}
\end{aligned}\end{equation}

\section{}
実対称行列の性質より、\(A\)は直交行列
\begin{equation}\begin{aligned}[b]
    P = \begin{pmatrix}
        v_1 & v_2 & v_3
    \end{pmatrix}
\end{aligned}\end{equation}
を用いて
\begin{equation}\begin{aligned}[b]
    D = P^\mathsf{T}AP = \text{diag}(1, 2, 4)
\end{aligned}\end{equation}
のように対角化できる。
よって
\begin{equation}\begin{aligned}[b]
    A^n &= \qty(PDP^\mathsf{T})^n\\
    &= PD^nP^\mathsf{T}\\
    &=v_1v_1^\mathsf{T}+2^nv_2v_2^\mathsf{T}+4^nv_3v_3^\mathsf{T}\\
    &=\frac{1}{6}\begin{pmatrix}
        2+3\cdot 2^n+2^{2n} & 2-2^{2n+1} & 2-3\cdot 2^n+2^{2n}\\
        2-2^{2n+1} & 2+2^{2n+2} & 2-2^{2n}+1\\
        2-3\cdot 2^n+2^{2n} & 2-2^{2n}+1 & 2+3\cdot 2^n+2^{2n}
    \end{pmatrix}
\end{aligned}\end{equation}

\section{}
\begin{equation}\begin{aligned}[b]
    (\lambda I-A)x &= b\\
    P^\mathsf{T}(\lambda I-A)P P ^\mathsf{T}x &= P^\mathsf{T} b\\
    x &= P(\lambda I-D)^{-1}P^\mathsf{T}b\\
    x^\mathsf{T}x &= b^\mathsf{T} P(\lambda I-D)^{-1}(\lambda I-D)^{-1}P^\mathsf{T}b\\
    &= \frac{\abs{v_1^\mathsf{T}b}^2}{(\lambda-1)^2} + \frac{\abs{v_2^\mathsf{T}b}^2}{(\lambda-2)^2} +\frac{\abs{v_3^\mathsf{T}b}^2}{(\lambda-4)^2}\\
    &= \frac{3}{(\lambda-1)^2} + \frac{2}{(\lambda-2)^2} +\frac{6}{(\lambda-4)^2}
\end{aligned}\end{equation}


\end{document}
